\section{Mode opératoire}

\subsection{Installation et Prérequis}
\noindent
Pour utiliser le compilateur \texttt{decac}, votre environnement de travail doit disposer des outils suivants :
\begin{itemize}[label=$\to$]
    \item Java (version 17 ou supérieure).
    \item Maven (pour la gestion des dépendances et la construction).
    \item La machine IMA installée et ajoutée dans votre \texttt{PATH}. 
\end{itemize}

\noindent
Avant votre première utilisation, compilez le projet à la racine :
\begin{command}
mvn package
\end{command}
\noindent
Le script \texttt{decac} est alors généré dans le répertoire \texttt{src/main/bin/}. Pour afficher toutes les options disponibles, exécutez simplement la commande \texttt{decac} sans arguments.

\subsection{Utilisation de l'extension Math}
\noindent
L'extension mathématique vous permet d'accéder aux fonctions trigonométriques précisées dans la partie précédente. Pour les utiliser, incluez le fichier d'en-tête dédié au début de votre fichier source \texttt{.deca} :
\begin{command}
#include "Math.decah"
\end{command}
\noindent
Le compilateur résout automatiquement l'inclusion et lie les méthodes nécessaires lors de la génération du code assembleur. Compilez votre fichier utilisant ces fonctions avec la commande :
\begin{command}
decac <fichier.deca>
\end{command}
\noindent
Puis, pour exécuter le code assembleur généré sur la machine IMA, utilisez :
\begin{command}
ima <fichier.ass>
\end{command}